% Generated draft: requires mathematical and source review.
\documentclass[11pt]{article}
\usepackage[margin=25mm]{geometry}
\usepackage{fontspec}
\setmainfont{Latin Modern Roman}
\usepackage{amsmath,amssymb,mathtools,unicode-math}
\setmathfont{Latin Modern Math}
\usepackage{enumitem,needspace,xurl,hyperref}
\setlength{\parindent}{0pt}
\setlength{\parskip}{5pt}
\setlength{\emergencystretch}{3em}
\newcommand{\N}{\mathbb{N}}
\newcommand{\Z}{\mathbb{Z}}
\newcommand{\Q}{\mathbb{Q}}
\newcommand{\R}{\mathbb{R}}
\newcommand{\C}{\mathbb{C}}
\newcommand{\F}{\mathbb{F}}
\newcommand{\A}{\mathbb{A}}
\newcommand{\PP}{\mathbb{P}}
\newcommand{\E}{\mathbb{E}}
\newcommand{\eps}{\varepsilon}
\newcommand{\dd}{\,\mathrm{d}}
\begin{document}
\section*{UnsolvedMath ID: 30002895}
\textbf{Source problem code:} catalogue entry
\subsection{Definitions and mathematical statement}
For notation only, if an evaluation of a matrix-valued function is a finite matrix
\[
A=F(X_1,\ldots,X_m)\in\mathbb K^{r\times s},\qquad \mathbb K\in\{\R,\C\},
\]
write \(q=\min\{r,s\}\). Its singular values \(\sigma_1(A)\ge\cdots\ge\sigma_q(A)\ge0\) are the square roots of the \(q\) largest eigenvalues of \(A^*A\), counted with multiplicity, where \(A^*\) denotes the conjugate transpose. Here ``small singular values'' refers to singular values near zero; the smallest singular value is \(\sigma_q(A)\). This notation does not impose a domain or matrix shape on the source problem.

\textbf{Source question.} Can the small singular values of general matrix-valued functions \(F(X_1,\ldots,X_m)\) be controlled?

\textbf{Unresolved formulation.} The supplied record does not specify the admissible functions \(F\), the domains or dimensions of the inputs and outputs, whether the inputs are random matrices or operators, or any probability laws, dependence assumptions, or limiting regime. It also does not specify what ``controlled'' means: no bound, probability estimate, integrability requirement, uniformity condition, or asymptotic rate is stated. No assumptions excluding identically singular evaluations are supplied. Consequently, the record supports this research question but does not determine a precise quantified conjecture. The reference to small singular values in the question and to smallest singular values in the supplied literature assessment does not resolve whether the intended target concerns only the minimum or a larger collection of values near zero.
\subsection{Short English statement}
The problem asks for control of singular values near zero for evaluations \(F(X_1,\ldots,X_m)\). The function class, input assumptions, and required estimates are unspecified, so a precise conjecture cannot be recovered from this record. Its open status is a supplied curation label requiring independent review.
\subsection{Sources}
\begin{itemize}
\item Local UnsolvedMath registry: permanent ID \(30002895\), title ``Small Singular Values of Matrix-Valued Free Functions,'' source problem code \texttt{catalogue entry}, in the Oberwolfach Reports Open Problems collection. The record lists proposed year 2015 and names no individual proposer.
\item Supplied source citation: \emph{Free Probability Theory (2016)}, Oberwolfach Reports, DOI \texttt{10.4171/owr/2015/28}, pp.~1571--1629. \url{https://doi.org/10.4171/owr/2015/28}. The background describes Oberwolfach Reports as documenting mathematical workshops and their open-problem discussions. The citation's 2016 label and the record's 2015 year information are retained without reconciliation.
\item Supplied literature source: Helton--Mai--Speicher, Journal of Functional Analysis. \url{https://doi.org/10.1016/j.jfa.2017.10.003}.
\item Supplied literature source: \emph{Asymptotic Spectra of Matrix-Valued Functions}. \url{https://arxiv.org/abs/1408.1732}.
\item Supplied literature source: Oberwolfach report 28. \url{https://doi.org/10.4171/owr/2015/28}.
\item Supplied literature assessment, checked \texttt{2026-08-21}: the record labels the problem open and states that free-probability linearization and realization methods solve limiting spectral distributions for broad polynomial and rational classes, while general control of the smallest singular values for arbitrary matrix-valued functions remains unresolved. This is supplied dated curation derived from the OWR material, not an independently verified literature assessment; its status and cited claims require independent verification.
\item The registry labels the statement \texttt{corrected\_verified}, explaining that only safe notation, typography, and TeX-delimiter cleanup was applied and that hypotheses, quantifiers, and subquestions were preserved. That supplied label does not establish that the formulation is mathematically complete or independently verified.
\end{itemize}
\end{document}
